\documentclass{article}
\usepackage{graphicx}
\usepackage{float}
\setlength{\parindent}{0pt}
\usepackage{caption}
\usepackage{fancyvrb}
\usepackage{tcolorbox}
\usepackage[a4paper, total={6in, 8.5in}]{geometry}
\setlength{\footskip}{80pt}
\usepackage{tabularx}
\usepackage{booktabs}
\usepackage{enumitem}

\title{Project - Security Hardware}
\author{Aleksander Kłap}
\date{08.06.2026}

\begin{document}

\maketitle
\pagebreak

\section{Security Hardware}
This chapter lists the hardware required to implement the redesigned architecture (Chapter 7). 

\subsection{Industrial Firewall}

\renewcommand{\arraystretch}{1.3} % 30% more row height
\begin{tabularx}{\textwidth}{>{\bfseries}p{2.5cm} X}
    \toprule
    Model & Fortinet FortiGate Rugged 70G \\
    \midrule
    Function & Central Layer 3 firewall and inter-VLAN router. All traffic between VLANs passes through it - under default-deny polict. 
    Can perform stateful filtering and packet inspection of industrial protocols (S7comm+, OPC UA, MQTT, HTTP/S). Implements four logical firewalls 
    (Edge, IT–IDMZ, IDMZ–OT Cell, OT Cell–OT Production) as separate rule sets between VLAN interfaces. \\
    \midrule
    Placement & Network core. Has one tagged VLAN interface (.254 gateway) per VLAN. Has single 802.1Q trunk to the managed switch. No VLAN is bridged on the switch, so the firewall is the only inter-VLAN path. \\
    \midrule
\end{tabularx}

\subsection{Managed Switch}
\renewcommand{\arraystretch}{1.3} % 30% more row height
\begin{tabularx}{\textwidth}{>{\bfseries}p{2.5cm} X}
    \toprule
    Model & Siemens SCALANCE XC-216  \\
    \midrule
    Function & VLAN-aware Layer 2 switching with IEEE 802.1Q tagging. Provides per-station access ports (untagged, one VLAN each) and the trunk uplink to the firewall. \\
    \midrule
    Placement & Replace the existing unmanaged SCALANCE XB008 switches at each station. Interconnected by 802.1Q trunks; one trunk reaches the firewall. \\
    \midrule
\end{tabularx}

\subsection{Engineering Station}
\renewcommand{\arraystretch}{1.3} % 30% more row height
\begin{tabularx}{\textwidth}{>{\bfseries}p{2.5cm} X}
    \toprule
    Model & Siemens SIMATIC IPC627E\\
    \midrule
    Function & Dedicated managed station running TIA Portal for PLC and HMI programming/configuration. Only host permitted S7comm+ to PLCs/HMIs.\\
    \midrule
    Placement & OT Cell zone, isolated in its own VLAN (40). Reaches OT Production only through the firewall on TCP 102.\\
    \midrule
\end{tabularx}

\subsection{IDMZ proxy server}
\renewcommand{\arraystretch}{1.3} % 30% more row height
\begin{tabularx}{\textwidth}{>{\bfseries}p{2.5cm} X}
    \toprule
    Model & Industrial rack server or SIMATIC IPC. Hardened Linux.\\
    \midrule
    Function & Hosts the IDMZ services — OPC UA proxy, Nginx reverse proxy, MQTT broker, NTP master, log forwarder, backup staging. \\
    \midrule
    Placement & Placed inside IDMZ zone - on VLAN20. \\
    \midrule
\end{tabularx}

\subsection{Syslog \& Backup Storage}
\renewcommand{\arraystretch}{1.3} % 30% more row height
\begin{tabularx}{\textwidth}{>{\bfseries}p{2.5cm} X}
    \toprule
    Model & Standard server, Linux OS. \\
    \midrule
    Function & Central Syslog collector for firewall, switches, proxies and broker; offline configuration backups (after every change) and daily MES database backups over SFTP.\\
    \midrule
    Placement & IT zone, VLAN10\\
    \midrule
\end{tabularx}

\end{document}

